\answerpaper{数学类}

\answerproblem{一、}{15}
设 $S$ 是空间中的一个椭球面．设方向为常向量 $V$ 的一束平行光照射 $S$，其中部分光线与 $S$ 相切，它们的切点在 $S$ 上形成一条曲线 $\Gamma$．证明：$\Gamma$ 落在一张过椭球中心的平面上．
\begin{solution}
\textbf{证明 1.} 在空间中取直角坐标系，记椭球面 $S$ 的方程为
\[
  \frac{x^2}{a^2}+\frac{y^2}{b^2}+\frac{z^2}{c^2}=1,\qquad V=(\alpha,\beta,\gamma).
\]
设 $(x,y,z)\in\Gamma$，则光束中的光线
\[
  l(t)=(x,y,z)+t(\alpha,\beta,\gamma),\qquad t\in\R
\]
是椭球面 $S$ 的切线．

由于每条切线与椭球面有且仅有一个交点，故 $t=0$ 是方程
\[
  \frac{(x+t\alpha)^2}{a^2}+\frac{(y+t\beta)^2}{b^2}+\frac{(z+t\gamma)^2}{c^2}=1
\]
的唯一解．由于 $(x,y,z)\in\Gamma\subset S$，上述方程化为
\[
  \left(\frac{\alpha^2}{a^2}+\frac{\beta^2}{b^2}+\frac{\gamma^2}{c^2}\right)t^2
  +2\left(\frac{\alpha}{a^2}x+\frac{\beta}{b^2}y+\frac{\gamma}{c^2}z\right)t=0.
\]
这个方程只有 $t=0$ 的唯一解，当且仅当
\[
  \frac{\alpha}{a^2}x+\frac{\beta}{b^2}y+\frac{\gamma}{c^2}z=0.
\]
这是一个过原点的平面方程，故 $\Gamma$ 落在过椭球中心的一张平面上．

\textbf{证明 2.} 在空间中做仿射变换，将椭球面映成圆球面．这时平行光束映成平行光束，切线映成切线，切点映成切点，椭球中心映成球面中心．

由于平行光束照圆球面的所有切线的切点是一个大圆，它落在过球心的平面上，而仿射变换将平面映成平面，故 $\Gamma$ 落在一张过椭球面中心的平面上．
\end{solution}

\answerproblem{二、}{15}
设 $n$ 为奇数，$A,B$ 为两个实 $n$ 阶方阵，且 $BA=0$．记 $A+J_A$ 的特征值集合为 $S_1$，$B+J_B$ 的特征值集合为 $S_2$，其中 $J_A,J_B$ 分别表示 $A,B$ 的 Jordan 标准型．求证：$0\in S_1\cup S_2$．
\begin{solution}
由秩不等式 $\rank A+\rank B\le\rank(BA)+n$，得 $\rank A+\rank B\le n$．结果
\[
  \rank A\le\frac{n}{2}\quad\text{或}\quad\rank B\le\frac{n}{2}.
\]
注意到 $n$ 为奇数，故有 $\rank A<\dfrac{n}{2}$ 或 $\rank B<\dfrac{n}{2}$ 成立．

若 $\rank A<\dfrac{n}{2}$，则 $\rank(A+J_A)\le\rank A+\rank J_A<n$，故 $0\in S_1$；

或 $\rank B<\dfrac{n}{2}$，则 $\rank(B+J_B)\le\rank B+\rank J_B<n$，故 $0\in S_2$．

所以最终有 $0\in S_1\cup S_2$．
\end{solution}

\answerproblem{三、}{20}
设 $A_1,\cdots,A_{2017}$ 为 $2016$ 阶实方阵．证明关于 $x_1,\cdots,x_{2017}$ 的方程
\[
  \det(x_1A_1+\cdots+x_{2017}A_{2017})=0
\]
至少有一组非零实数解，其中 $\det$ 表示行列式．
\begin{solution}
记
\[
  A_1=\left(p_1^{(1)},\cdots,p_{2016}^{(1)}\right),\qquad\cdots,\qquad
  A_{2017}=\left(p_1^{(2017)},\cdots,p_{2016}^{(2017)}\right).
\]
考虑线性方程组
\[
  x_1p_1^{(1)}+\cdots+x_{2017}p_1^{(2017)}=\mathbf{0}.
\]
由于未知数个数大于方程个数，故该线性方程组必有非零解 $(c_1,\cdots,c_{2017})$．从而 $c_1A_1+\cdots+c_{2017}A_{2017}$ 的第一列为 $\mathbf{0}$，更有
\[
  \det(c_1A_1+\cdots+c_{2017}A_{2017})=0.
\]
\end{solution}

\answerproblem{四、}{20}
设 $f_0(x),f_1(x)$ 是 $[0,1]$ 上正连续函数，满足 $\displaystyle\int_0^1 f_0(x)\dd x\le\int_0^1 f_1(x)\dd x$．

设
\[
  f_{n+1}(x)=\frac{2f_n^2(x)}{f_n(x)+f_{n-1}(x)},\qquad n=1,2,\cdots.
\]
求证：数列 $\displaystyle a_n=\int_0^1 f_n(x)\dd x$，$n=1,2,\cdots$ 单调递增且收敛．
\begin{solution}
因为
\begin{align*}
  &\int_0^1\frac{f_1^2(x)}{f_1(x)+f_0(x)}\dd x-\int_0^1\frac{f_0^2(x)}{f_1(x)+f_0(x)}\dd x\\
  &\quad=\int_0^1\frac{f_1^2(x)-f_0^2(x)}{f_1(x)+f_0(x)}\dd x
  =\int_0^1 f_1(x)\dd x-\int_0^1 f_0(x)\dd x\ge 0.
\end{align*}
所以
\begin{align*}
  a_2-a_1&=2\int_0^1\frac{f_1^2(x)}{f_1(x)+f_0(x)}\dd x-\int_0^1 f_1(x)\dd x\\
  &=\int_0^1\frac{f_1^2(x)}{f_1(x)+f_0(x)}\dd x-\int_0^1\frac{f_1(x)f_0(x)}{f_1(x)+f_0(x)}\dd x\\
  &\ge\frac{1}{2}\int_0^1\frac{f_1^2(x)+f_0^2(x)}{f_1(x)+f_0(x)}\dd x-\int_0^1\frac{f_1(x)f_0(x)}{f_1(x)+f_0(x)}\dd x\\
  &=\int_0^1\frac{\left[f_1(x)-f_0(x)\right]^2}{2\left[f_1(x)+f_0(x)\right]}\dd x\ge 0.
\end{align*}
归纳地可以证明 $a_{n+1}\ge a_n$，$n=1,2,\cdots$．

由于 $f_0,f_1$ 为正的连续函数，可取常数 $k\ge 1$，使得 $f_1\le kf_0$．设 $c_1=k$．根据递推关系式可以归纳证明
\begin{equation}
  f_n(x)\le c_nf_{n-1}(x),\tag{1}
\end{equation}
其中 $c_{n+1}=\dfrac{2c_n}{c_n+1}$，$n=0,1,\cdots$．容易证明 $\{c_n\}$ 单调递减且趋于 $1$，且
\[
  \frac{c_n}{c_n+1}\le\frac{k}{k+1}.
\]
以下证明 $\{a_n\}$ 收敛．由 (1) 可得 $a_{n+1}\le c_{n+1}a_n$．因此
\[
  c_{n+1}a_{n+1}\le\frac{2c_{n+1}}{c_n+1}c_na_n
  =\frac{4c_n}{(c_n+1)^2}c_na_n\le c_na_n.
\]
这就说明 $\{c_na_n\}$ 是正单调递减数列，因而收敛．注意到 $\{c_n\}$ 收敛到 $1$，可知 $\{a_n\}$ 收敛，且有 $\lim\limits_{n\to\infty}a_n\le c_1a_1=ka_1$．
\end{solution}

\answerproblem{五、}{15}
设 $\alpha>1$．求证：不存在 $[0,+\infty)$ 上的正可导函数 $f(x)$ 满足
\begin{equation}
  f'(x)\ge f^{\alpha}(x),\qquad x\in[0,+\infty).\tag{1}
\end{equation}
\begin{solution}
若 $f(x)$ 是这样的函数，则 $f'(x)>0$．因此 $f(x)$ 是严格递增函数．(1) 可以表示为
\[
  \left(\frac{1}{\alpha-1}f^{1-\alpha}(x)+x\right)'\le 0.
\]
这说明 $\dfrac{1}{\alpha-1}f^{1-\alpha}(x)+x$ 是单调递减函数．因而
\[
  \frac{1}{\alpha-1}f^{1-\alpha}(x+1)+(x+1)\le\frac{1}{\alpha-1}f^{1-\alpha}(x)+x,
\]
即 $\alpha-1\le f^{1-\alpha}(x)-f^{1-\alpha}(x+1)<f^{1-\alpha}(x)$．因此有 $f^{\alpha-1}(x)<\dfrac{1}{\alpha-1}$．从而 $f(x)$ 是有界函数．

从 $f(x)$ 的严格递增性，可知 $\lim\limits_{x\to+\infty}f(x)$ 收敛．由微分中值定理，存在 $\xi\in(x,x+1)$，使得
\[
  f(x+1)-f(x)=f'(\xi)\ge f^{\alpha}(x)\ge f^{\alpha}(0)>0.
\]
令 $x\to+\infty$，上式左端趋于 $0$，可得矛盾！
\end{solution}

\answerproblem{六、}{15}
设 $f(x),g(x)$ 是 $[0,1]$ 区间上的单调递增函数，满足
\[
  0\le f(x),g(x)\le 1,\qquad \int_0^1 f(x)\dd x=\int_0^1 g(x)\dd x.
\]
求证：$\displaystyle\int_0^1\abs{f(x)-g(x)}\dd x\le\frac{1}{2}$．
\begin{solution}
由于 $f,g$ 可用单调阶梯函数逼近，故可不妨设它们都是单调增的阶梯函数．令 $h(x)=f(x)-g(x)$，则对 $\forall x,y\in[0,1]$，有 $\abs{h(x)-h(y)}\le 1$．

事实上，对 $x\ge y$，我们有
\begin{align*}
  -1\le -\bigl(g(x)-g(y)\bigr)\le{}&h(x)-h(y)\\
  &{}=f(x)-f(y)-\bigl(g(x)-g(y)\bigr)\le f(x)-f(y)\le 1.
\end{align*}
对 $x<y$，有
\[
  -1\le f(x)-f(y)\le h(x)-h(y)\le g(y)-g(x)\le 1.
\]
现记
\[
  C_1=\{\,x\in[0,1]\mid f(x)\ge g(x)\,\},\qquad
  C_2=\{\,x\in[0,1]\mid f(x)<g(x)\,\},
\]
则 $C_1,C_2$ 分别为有限个互不相交区间的并，且由 $\displaystyle\int_0^1 f\dd x=\int_0^1 g\dd x$，有
\[
  \int_{C_1}h\dd x=-\int_{C_2}h\dd x.
\]
让 $\abs{C_i}\ (i=1,2)$ 表示 $C_i$ 所含的那些区间的长度之和，则
\[
  \abs{C_1}+\abs{C_2}=1.
\]
于是
\begin{align*}
  2\int_0^1\abs{f-g}\dd x
  &=2\left(\int_{C_1}h\dd x-\int_{C_2}h\dd x\right)\\
  &\le\left(\frac{\abs{C_2}}{\abs{C_1}}\int_{C_1}h\dd x
    +\frac{\abs{C_1}}{\abs{C_2}}\int_{C_2}(-h)\dd x\right)
    +\int_{C_1}h\dd x-\int_{C_2}h\dd x\\
  &=\left(\frac{1}{\abs{C_1}}\int_{C_1}h\dd x
    +\frac{1}{\abs{C_2}}\int_{C_2}(-h)\dd x\right)\\
  &\le\sup_{C_1}h+\sup_{C_2}(-h)\le 1.
\end{align*}
注意，上式中最后一个不等式来自 $\abs{h(x)-h(y)}\le 1$，另外，若有某个 $\abs{C_i}$ 等于 $0$，则结论显然成立．
\end{solution}
